\section{Verus}

\begin{frame}{Verus connects Rust code to Z3}
\begin{center}
\begin{tikzpicture}[node distance=6mm,box/.style={draw=Teal,rounded corners,align=center,minimum height=15mm,text width=28mm,font=\small}]
\node[box] (s) {Rust code\\+ contracts};
\node[box,right=of s] (v) {Verus\\obligations};
\node[box,right=of v] (z) {Z3\\checks};
\draw[-{Stealth},thick,Teal] (s)--(v);
\draw[-{Stealth},thick,Teal] (v)--(z);
\end{tikzpicture}
\end{center}
\begin{itemize}
\item Write specifications and proofs alongside executable code.
\item Verify supported Rust code with annotations and library specifications.
\item Erase specification/proof code when compiling the executable.
\end{itemize}
\source{https://github.com/verus-lang/verus}{Verus project}\source{https://verus-lang.github.io/verus/guide/modes.html}{Verus Guide: code modes}
\note{Verus supports a subset of Rust; existing crates are not automatically fully verified. The diagram is conceptual, not a complete description of the internal compiler pipeline.}
\end{frame}

\begin{frame}{Three modes, three jobs}
\begin{tabular}{p{.17\textwidth}p{.38\textwidth}p{.32\textwidth}}
\toprule
Mode & Purpose & Example \\\midrule
\texttt{exec} & Actual computation & Compute a maximum \\\addlinespace
\texttt{spec} & Mathematical definitions & Describe the desired result \\\addlinespace
\texttt{proof} & Establish logical facts & Supply a lemma or assertion \\\bottomrule
\end{tabular}
\vspace{5mm}
\begin{itemize}
\item Think: computation, definition, and lemma.
\item Specifications can use mathematical \texttt{int} and \texttt{nat}.
\item Ghost code is reasoning-only code, erased before execution.
\end{itemize}
\source{https://verus-lang.github.io/verus/guide/modes.html}{Verus Guide: code modes}
\note{Machine integers still matter in exec code. Mathematical types let the specification describe arithmetic without silently wrapping. Proof functions can package reusable arguments.}
\end{frame}

\begin{frame}[fragile]{Implementing max in Rust}
\begin{lstlisting}
fn max(x: i64, y: i64) -> i64
{
    if x>y {x} else {y}
}
\end{lstlisting}
\end{frame}
\begin{frame}[fragile]{Decorating our implementation with Verus}
\begin{lstlisting}
fn max(x: i64, y: i64) -> (m: i64)
    requires
    ensures
        m>=x,
        m>=y,
        (m==x) || (m==y)
{
    if x>y {x} else {y}
}
\end{lstlisting}
\begin{itemize}
\item \texttt{requires}: what callers must establish.
\item \texttt{ensures}: what callers may rely on afterward.
\end{itemize}
\source{https://verus-lang.github.io/verus/guide/requires_ensures.html}{Verus Guide: contracts}
\note{The complete file is examples/max.rs. There are no preconditions on the two i64 inputs. The executable function only compares and selects, so there is no arithmetic-overflow obligation for its body. The three postconditions characterize the maximum.}
\end{frame}

\begin{frame}{Using a verified function is applying a lemma}
Let $M(a,b)$ denote the verified function. Its contract gives
\[
 M(a,b)\ge a,\qquad M(a,b)\ge b,\qquad M(a,b)\in\{a,b\}.
\]
To find the maximum of three inputs, compute $t=M(a,b)$, then $m=M(t,c)$.
\begin{itemize}
\item $m\ge t\ge a,b$ and $m\ge c$: $m$ is an upper bound.
\item $m\in\{t,c\}$ and $t\in\{a,b\}$: $m\in\{a,b,c\}$.
\end{itemize}
\begin{block}{Modular proof}
The two contracts establish $m=\max\{a,b,c\}$.\par
We need not inspect either function call's implementation.
\end{block}
\note{All inputs are representable machine integers; comparison and selection stay in the same type. This is ordinary lemma application followed by transitivity. Repeating this argument suggests the maximum-over-a-slice algorithm on the next slides.}
\end{frame}

\begin{frame}[fragile]{Maximum over a nonempty \texttt{u64} slice}
\lstinputlisting[firstline=5,lastline=9]{examples/max_iter.rs}
\begin{itemize}
\item Nonempty input: a maximum must be attained by some element.
\item First postcondition: $m$ bounds every element from above.
\item Second postcondition: an element actually equals $m$.
\end{itemize}
\note{The binary example used i64. For this loop example we use u64 as requested; the same order argument works for both. A slice is a sequence viewed without copying it. The mathematical int indices in the quantifiers are ghost expressions, not runtime iteration.}
\end{frame}

\begin{frame}{What about loops?}
    Verus handles loops in rust by induction.
    
    For each loop, we must establish that
    \begin{itemize}
        \item it terminates
        \item the induction hypothesis implies the next step
    \end{itemize}
    \begin{block}{Proof idea}
        \small
        We assume \texttt{values} is non-empty, so use induction on the length of \texttt{values}.
        \begin{itemize}
            \setlength{\itemsep}{3pt}
            \item If \texttt{values.len() == 1}, then the max is \texttt{values[0]}.
            \item For the induction hypothesis, assume that for all
                $k < \texttt{values.len()}$, $\texttt{values[0..k]} \le \texttt{values[j]}$, for some $j<k$.
            \item We can then consider \texttt{values[k+1]}---if it's greater than the current max use it, otherwise stick with what we have.
            \item This process terminates because \texttt{values} is of finite length.
        \end{itemize}
    \end{block}
\end{frame}

\begin{frame}[fragile]{The induction hypothesis in Verus}
\lstinputlisting[firstline=10,lastline=26,basicstyle=\ttfamily\fontsize{8}{9}\selectfont]{examples/max_iter.rs}
\source{https://verus-lang.github.io/verus/guide/while.html}{Verus Guide: Loops and invariants}
\end{frame}

\begin{frame}{What is being trusted?}
\begin{itemize}
\item The specification captures the intended theorem.
\pause
\item The translation from code to logical obligations is sound.
\pause
\item The solver's result and any assumed lemmas are correct.
\pause
\item Compilation and execution preserve the verified meaning.
\end{itemize}
\end{frame}

